\scanpage{179}
\mspage{95}
% archive sheet number: 178
% line 1, continuation
donc le lemme 4.4.10 ne peut être invoqué.

% lines 2-12
2. Nous venons de prouver, via 4.4.8 et 4.4.9 (dont la démonstration est achevée avec 4.4.12), que \deleted{\unsure{le}} si $C$ est une petite catégorie stable par $\unsure{\varinjlim|\underline{\lim}}$ de cardinal $\leq\pi$, alors $\operatorname{Ind}_{\pi}C$ est également stable par ces limites. \deleted{Donc} \inserted{De plus} \deleted{pour} les \deleted{limites de type $J$} systèmes inductifs de type $J$ dans $\operatorname{Ind}_{\pi}C$ \deleted{\unsure{rep|resp.|se}} admettent une représentation indicielle
\[
\deleted{\unsure{\varphi}}\quad\psi:I\times J\mathrel{\rule{1.5em}{0.4pt}}C,
\]
avec $I$ grand devant $\pi$. Et on peut toujours supposer, quitte à prendre $I'\longrightarrow I$ cofinal avec $I'$ ens. ordonné grand devant $\pi$, que $I$ est \inserted{de plus} un ens. ordonné.

% lines 18-20
C'est là le \unsure{noyau} \unsure{dur|du}, pour la démonstration du

% lines 21-25
\uline{Théorème 4.4.13.} Soit $C$ une catégorie équivalente à une petite catégorie, et stable par $\varinjlim$ de cardinal $\leq\pi$, où $\pi$ est un cardinal infini donné. Alors: \deleted{$\operatorname{Ind}_{\pi}$ \unsure{jouit|pour} \illegible{1}}
